% GENERATED FILE. Edit canonical lessons, then run npm run generate:books.
% canonical-source-hash: fnv1a64:dd295a07a84622a0
% canonical-lessons: JA-C345-henkou-suru, JA-C345-enki-suru, JA-C345-torikesu, JA-C345-zurasu, JA-C345-hayameru

\chapter{When Plans Change}
\label{ch:ja-when-plans-change}
\hlchaptermodality{345}

\begin{chapteropening}
\textbf{By the end of this chapter:} \emph{I can change, put off or bring forward a plan if something happens.}

Complete the last lesson of chapter 345: I can change, put off or bring forward a plan if something happens.
\end{chapteropening}

\section[\texorpdfstring{\ja{へんこうする} (henkō suru)}{へんこうする (henkō suru)}]{\ja{へんこうする} (henkō suru) — to change a plan; to alter (a date)}

\label{lesson:JA-C345-henkou-suru}



\begin{quote}
Before the new one: say the Japanese for a schedule, then the Japanese for going smoothly.
\end{quote}

\subsection*{You'll want to know: \ja{へんこうする}}
\textbf{\ja{へんこうする}} — \emph{henkō suru} — \textquotedblleft{}to change a plan; to alter (a date)\textquotedblright{}.

\begin{quote}
\textbf{\ja{もし} \ja{あめが} \ja{ふったら、よていを} \ja{へんこうします。}} — \emph{moshi ame ga futtara, yotei o henkō shimasu} — \textquotedblleft{}If it rains, we will change the plan.\textquotedblright{}
\end{quote}

Polite \textbf{\ja{へんこうします}}, \emph{henkōshimasu}; a \ja{する} verb. Past (polite): \textbf{\ja{かいしゃが} \ja{よていを} \ja{へんこうしました。}}, \emph{kaisha ga yotei o henkō shimashita}.

\begin{cousinweb}
\textbf{\ja{へんこうする}} (\emph{henkō suru}) is built from a Sino-Japanese noun and the native verb \emph{suru}, \textquotedblleft{}to do\textquotedblright{}. The noun, \emph{henkō}, is two Chinese characters read with their on'yomi, readings borrowed from Chinese: \emph{hen}, \textquotedblleft{}change\textquotedblright{} (Middle Chinese \emph{pjenH}) and \emph{kō}, \textquotedblleft{}renew\textquotedblright{} (\emph{kaeng}) (KANJIDIC2; JMdict; Middle Chinese in Baxter's transcription).
\end{cousinweb}

\begin{grammarlens}[title={What you've built}]
One more word for saying what happens if.
\end{grammarlens}

\subsection*{Your turn}
\begin{itemize}
\raggedright
  \item \emph{Say it:} \emph{henkō suru}
  \item \emph{Say it:} \emph{henkō suru}, once more
  \item \emph{Say it:} \emph{moshi ame ga futtara, yotei o henkō shimasu} — If it rains, we will change the plan.
  \item \emph{Recall:} say the Japanese for a schedule, then the Japanese for going smoothly, then say \emph{henkō suru} again
\end{itemize}

\begin{tcolorbox}[breakable,colback=teal!4,colframe=teal!35!black,title={Before you move on}]
What does \ja{へんこうする} mean? (\textquotedblleft{}To change a plan; to alter (a date)\textquotedblright{}.) Say it once more.
\end{tcolorbox}
\section[\texorpdfstring{\ja{えんきする} (enki suru)}{えんきする (enki suru)}]{\ja{えんきする} (enki suru) — to postpone}

\label{lesson:JA-C345-enki-suru}



\begin{quote}
Before the new one: say the Japanese for going smoothly, then the Japanese for to change a plan.
\end{quote}

\subsection*{You'll want to know: \ja{えんきする}}
\textbf{\ja{えんきする}} — \emph{enki suru} — \textquotedblleft{}to postpone\textquotedblright{}.

\begin{quote}
\textbf{\ja{たいふうが} \ja{きたら、しあいを} \ja{えんきします。}} — \emph{taifū ga kitara, shiai o enki shimasu} — \textquotedblleft{}If the typhoon comes, we will postpone the match.\textquotedblright{}
\end{quote}

Polite \textbf{\ja{えんきします}}, \emph{enkishimasu}; a \ja{する} verb. Past (polite): \textbf{\ja{がっこうが} \ja{しあいを} \ja{えんきしました。}}, \emph{gakkō ga shiai o enki shimashita}.

\begin{cousinweb}
\textbf{\ja{えんきする}} (\emph{enki suru}) is built from a Sino-Japanese noun and the native verb \emph{suru}, \textquotedblleft{}to do\textquotedblright{}. The noun, \emph{enki}, is two Chinese characters read with their on'yomi, readings borrowed from Chinese: \emph{en}, \textquotedblleft{}prolong\textquotedblright{} (Middle Chinese \emph{yen}) and \emph{ki}, \textquotedblleft{}period\textquotedblright{} (\emph{gi}) (KANJIDIC2; JMdict; Middle Chinese in Baxter's transcription).
\end{cousinweb}

\begin{grammarlens}[title={What you've built}]
One more word for saying what happens if.
\end{grammarlens}

\subsection*{Your turn}
\begin{itemize}
\raggedright
  \item \emph{Say it:} \emph{enki suru}
  \item \emph{Say it:} \emph{enki suru}, once more
  \item \emph{Say it:} \emph{taifū ga kitara, shiai o enki shimasu} — If the typhoon comes, we will postpone the match.
  \item \emph{Recall:} say the Japanese for going smoothly, then the Japanese for to change a plan, then say \emph{enki suru} again
\end{itemize}

\begin{tcolorbox}[breakable,colback=teal!4,colframe=teal!35!black,title={Before you move on}]
What does \ja{えんきする} mean? (\textquotedblleft{}To postpone\textquotedblright{}.) Say it once more.
\end{tcolorbox}
\section[\texorpdfstring{\ja{とりけす} (torikesu)}{とりけす (torikesu)}]{\ja{とりけす} (torikesu) — to cancel; to take back}

\label{lesson:JA-C345-torikesu}



\begin{quote}
Before the new one: say the Japanese for to change a plan, then the Japanese for to postpone.
\end{quote}

\subsection*{You'll want to know: \ja{とりけす}}
\textbf{\ja{とりけす}} — \emph{torikesu} — \textquotedblleft{}to cancel; to take back\textquotedblright{}.

\begin{quote}
\textbf{\ja{びょうきに} \ja{なったら、よやくを} \ja{とりけして} \ja{ください。}} — \emph{byōki ni nattara, yoyaku o torikeshite kudasai} — \textquotedblleft{}If you fall ill, please cancel the booking.\textquotedblright{}
\end{quote}

Polite \textbf{\ja{とりけします}}, \emph{torikeshimasu}; a godan verb. Past (polite): \textbf{\ja{ちちが} \ja{よやくを} \ja{とりけしました。}}, \emph{chichi ga yoyaku o torikeshimashita}.

\begin{cousinweb}
\textbf{\ja{とりけす}} (\emph{torikesu}) is a native Japanese word (wago), inherited rather than borrowed from Chinese: it is built from \emph{tori}, \textquotedblleft{}taking (toru)\textquotedblright{} and \emph{kesu}, \textquotedblleft{}to put out\textquotedblright{}, and its kanji were attached later as kun'yomi (JMdict).
\end{cousinweb}

\begin{grammarlens}[title={What you've built}]
One more word for saying what happens if.
\end{grammarlens}

\subsection*{Your turn}
\begin{itemize}
\raggedright
  \item \emph{Say it:} \emph{torikesu}
  \item \emph{Say it:} \emph{torikesu}, once more
  \item \emph{Say it:} \emph{byōki ni nattara, yoyaku o torikeshite kudasai} — If you fall ill, please cancel the booking.
  \item \emph{Recall:} say the Japanese for to change a plan, then the Japanese for to postpone, then say \emph{torikesu} again
\end{itemize}

\begin{tcolorbox}[breakable,colback=teal!4,colframe=teal!35!black,title={Before you move on}]
What does \ja{とりけす} mean? (\textquotedblleft{}To cancel; to take back\textquotedblright{}.) Say it once more.
\end{tcolorbox}
\section[\texorpdfstring{\ja{ずらす} (zurasu)}{ずらす (zurasu)}]{\ja{ずらす} (zurasu) — to shift; to move (a time) a little}

\label{lesson:JA-C345-zurasu}



\begin{quote}
Before the new one: say the Japanese for to postpone, then the Japanese for to cancel.
\end{quote}

\subsection*{You'll want to know: \ja{ずらす}}
\textbf{\ja{ずらす}} — \emph{zurasu} — \textquotedblleft{}to shift; to move (a time) a little\textquotedblright{}.

\begin{quote}
\textbf{\ja{かいぎが} \ja{ながく} \ja{なったら、ひるごはんを} \ja{ずらします。}} — \emph{kaigi ga nagaku nattara, hirugohan o zurashimasu} — \textquotedblleft{}If the meeting runs long, we will move lunch back a little.\textquotedblright{}
\end{quote}

Polite \textbf{\ja{ずらします}}, \emph{zurashimasu}; a godan verb. Past (polite): \textbf{\ja{せんせいが} \ja{じかんを} \ja{ずらしました。}}, \emph{sensei ga jikan o zurashimashita}.

\begin{cousinweb}
\textbf{\ja{ずらす}} (\emph{zurasu}) is a native Japanese word (wago), inherited rather than borrowed from Chinese, and its older origin is not settled; JMdict gives it no kanji spelling (JMdict).
\end{cousinweb}

\begin{grammarlens}[title={What you've built}]
One more word for saying what happens if.
\end{grammarlens}

\subsection*{Your turn}
\begin{itemize}
\raggedright
  \item \emph{Say it:} \emph{zurasu}
  \item \emph{Say it:} \emph{zurasu}, once more
  \item \emph{Say it:} \emph{kaigi ga nagaku nattara, hirugohan o zurashimasu} — If the meeting runs long, we will move lunch back a little.
  \item \emph{Recall:} say the Japanese for to postpone, then the Japanese for to cancel, then say \emph{zurasu} again
\end{itemize}

\begin{tcolorbox}[breakable,colback=teal!4,colframe=teal!35!black,title={Before you move on}]
What does \ja{ずらす} mean? (\textquotedblleft{}To shift; to move (a time) a little\textquotedblright{}.) Say it once more.
\end{tcolorbox}
\section[\texorpdfstring{\ja{はやめる} (hayameru)}{はやめる (hayameru)}]{\ja{はやめる} (hayameru) — to bring forward; to make earlier}

\label{lesson:JA-C345-hayameru}



\begin{quote}
Before the new one: say the Japanese for to cancel, then the Japanese for to shift a time.
\end{quote}

\subsection*{You'll want to know: \ja{はやめる}}
\textbf{\ja{はやめる}} — \emph{hayameru} — \textquotedblleft{}to bring forward; to make earlier\textquotedblright{}.

\begin{quote}
\textbf{\ja{ゆきが} \ja{ふったら、しゅっぱつを} \ja{はやめます。}} — \emph{yuki ga futtara, shuppatsu o hayamemasu} — \textquotedblleft{}If it snows, we will set off earlier.\textquotedblright{}
\end{quote}

Polite \textbf{\ja{はやめます}}, \emph{hayamemasu}; an ichidan verb. Past (polite): \textbf{\ja{ちちが} \ja{しゅっぱつを} \ja{はやめました。}}, \emph{chichi ga shuppatsu o hayamemashita}.

\begin{cousinweb}
\textbf{\ja{はやめる}} (\emph{hayameru}) is a native Japanese word (wago), inherited rather than borrowed from Chinese, and its older origin is not settled. Its kanji was attached to it later as a kun'yomi, a native reading given to a Chinese character (JMdict, KANJIDIC2).
\end{cousinweb}

\begin{grammarlens}[title={What you've built}]
One more word for saying what happens if.
\end{grammarlens}

\subsection*{Your turn}
\begin{itemize}
\raggedright
  \item \emph{Say it:} \emph{hayameru}
  \item \emph{Say it:} \emph{hayameru}, once more
  \item \emph{Say it:} \emph{yuki ga futtara, shuppatsu o hayamemasu} — If it snows, we will set off earlier.
  \item \emph{Recall:} say the Japanese for to cancel, then the Japanese for to shift a time, then say \emph{hayameru} again
\end{itemize}

\begin{tcolorbox}[breakable,colback=teal!4,colframe=teal!35!black,title={Before you move on}]
What does \ja{はやめる} mean? (\textquotedblleft{}To bring forward; to make earlier\textquotedblright{}.) Say it once more.
\end{tcolorbox}
